\begin{homeworkProblem}
Sea $x\in\mathbb{R}$. Es bien conocido (y puede demostrarse usando el
principio del palomar) que para todo irracional $x$ existen infinitos
racionales $p/q$ en forma reducida tales que
\begin{align*}
  \left|x-\frac{p}{q}\right|\le \frac{1}{q^{2}}. 
\end{align*}
En este ejercicio queremos probar que si se reemplaza $1/q^{2}$ por
$1/q^{3}$ (o más generalmente por $1/q^{2+\varepsilon}$), entonces el
conjunto de $x$ que admiten infinitas aproximaciones de ese tipo es de
medida cero.
\end{homeworkProblem}
\begin{solucion}
Para cada par $(p,q)$ con $(p,q)=1$, definamos el intervalo
\[
I_{p,q}
=
\left\{\,x\in\mathbb{R}:\,
\left|x-\frac{p}{q}\right|\le \frac{1}{q^{3}}
\right\}
=
\left[\frac{p}{q}-\frac{1}{q^{3}},\, \frac{p}{q}+\frac{1}{q^{3}}\right].
\]

Cada intervalo tiene longitud
\[
|I_{p,q}|=\frac{2}{q^{3}}.
\]

Para un denominador fijo $q$, consideremos todos los $p$ con $(p,q)=1$.
El número de tales $p$ está dado por la función totiente $\varphi(q)$, y
satisfacemos siempre la cota
\[
\varphi(q)\le q.
\]

Sea ahora
\[
E_q := \bigcup_{\substack{p\in\mathbb{Z}\\ (p,q)=1}} I_{p,q}.
\]

Entonces podemos acotar la medida de $E_q$ por
\[
m(E_q)\le \sum_{\substack{p\in\mathbb{Z}\\ (p,q)=1}} |I_{p,q}|
= \varphi(q)\frac{2}{q^{3}}
\le \frac{2}{q^{2}}.
\]

El conjunto de números $x$ que admiten infinitas aproximaciones del tipo
$\lvert x-\frac{p}{q}\rvert\le\frac{1}{q^{3}}$ es entonces
\[
E := \limsup_{q\to\infty} E_q
=
\bigcap_{N=1}^{\infty}\ \bigcup_{q\ge N} E_q.
\]

Aplicamos ahora el lema de Borel--Cantelli. Tenemos
\[
\sum_{q=1}^{\infty} m(E_q)
\le
\sum_{q=1}^{\infty} \frac{2}{q^{2}}
<\infty.
\]

Por el lema de Borel--Cantelli se deduce que
\[
m(E)=0.
\]

Por tanto, el conjunto de números reales $x$ para los cuales existen
infinitos racionales $p/q$ en forma reducida tales que
\[
\left|x-\frac{p}{q}\right|\le\frac{1}{q^{3}}
\]
tiene medida cero.

\medskip
\noindent\textbf{Extensión a $1/q^{2+\varepsilon}$.}
Con el mismo argumento, si definimos
\[
I_{p,q}
=
\left\{
x:\,
\left|x-\frac{p}{q}\right|\le \frac{1}{q^{2+\varepsilon}}
\right\},
\]
su longitud es $2/q^{2+\varepsilon}$ y se obtiene
\[
m(E_q)\le \varphi(q)\frac{2}{q^{2+\varepsilon}}
\le \frac{2}{q^{1+\varepsilon}}.
\]
La serie
\[
\sum_{q=1}^{\infty} \frac{1}{q^{1+\varepsilon}}
\]
converge para todo $\varepsilon>0$, de modo que el lema de
Borel--Cantelli vuelve a aplicar y el conjunto correspondiente es de
medida cero.

\end{solucion}
